\chapter*{Úvod}

Současné doporučovací systémy dosahují vysoké přesnosti díky komplexním latentním reprezentacím, které jsou však často neprůhledné a obtížně ovladatelné. Cílem této práce je využít řídké autoenkodéry (\textit{sparse autoencoders}) integrované do architektury kolaborativního filtrování k transformaci latentních vektorů na interpretovatelné a uživatelsky ovlivnitelné reprezentace preferencí v doméně volnočasových aktivit (tzv. \textit{points of interest}, POIs). Práce navazuje na aktuálně recenzovaný článek vedoucího práce \cite{Spisak2024} a adaptuje jeho výsledky do cílové domény doporučování míst.

Tento přístup reaguje na fundamentální problém moderních doporučovacích systémů. Přestože hluboké modely a maticové faktorizace přinášejí vynikající prediktivní přesnost, během trénování vytvářejí husté, propletené (\textit{entangled}) reprezentace, jež fungují jako tzv. černé skříňky (\textit{black-boxes}). Neprůhlednost těchto latentních reprezentací představuje zásadní překážku -- uživatelé nerozumí tomu, proč jim jsou konkrétní položky nabízeny, a vývojáři mají omezené možnosti odhalit interní mechanismy algoritmu \cite{Chen2022}. Tento nedostatek transparentnosti nejenže snižuje celkovou důvěru v systém \cite{Tintarev2022}, ale především omezuje uživatelskou ovladatelnost. Pokud uživatel nechápe profil, který si o něm systém vytvořil, nemá ani možnost jeho nepřesnosti korigovat a směřovat tak další doporučení správným směrem \cite{Ooge2023}.

\section*{Motivace a cíl práce}

V kontextu doporučování bodů zájmu (POI) je absence transparentnosti a interaktivity ještě kritičtější než v tradičních doménách (jako jsou e-shopy či streamovací platformy pro filmy). Preference v oblasti volnočasových aktivit jsou totiž silně vázány na geografické a časoprostorové faktory \cite{Zeng2025}, jako je fyzická vzdálenost, dostupnost či aktuální situace uživatele. Uživatelské záměry se v této doméně mění velice dynamicky. Běžné doporučovací modely nedokážou tento složitý kontext uživateli dostatečně srozumitelně zprostředkovat a neposkytují mu nástroje pro snadnou interaktivní úpravu v reálném čase.

Cílem této práce je proto vytvořit model využívající řídké autoenkodéry (SAE) pro extrakci vysvětlitelných a upravitelných preferencí. Začleněním mechanismu řídkosti (sparsity) do procesu kolaborativního filtrování se model naučí transformovat abstraktní doporučovací faktory do sémanticky pochopitelné podoby. Tyto extrahované rysy pak poslouží nejen jako základ pro generování textových či vizuálních vysvětlení, ale především jako interaktivní "ovládací prvky" (\textit{knobs}), skrze které může uživatel své preference přímo řídit a upravovat.

\section*{Přínos práce}
Hlavní přínos předkládané diplomové práce spočívá v inovativní syntéze pokročilých přístupů z oblasti interpretability do specifické domény doporučování bodů zájmu (POI). Tradiční modely kolaborativního filtrování sice dosahují vysoké přesnosti, ale vytvářejí neprůhledné latentní reprezentace, což je v doméně POI – navíc často zatížené extrémní řídkostí dat a nedostatkem detailních sémantických informací – zásadní překážka. Práce na tento stav reaguje návrhem architektury, která navazuje na moderní koncepty rozpletených (\textit{disentangled}) reprezentací \cite{Wang2021} a využívá řídké autoenkodéry (SAE) \cite{Spisak2024} k transformaci komplexních uživatelských preferencí do izolovaných rysů. K překonání absence sémantiky v lokačních datech je nově navržena integrace velkých jazykových modelů (LLM), které z historie návštěv uživatele automaticky extrahují tzv. pseudo-profily \cite{Zeng2025} a přiřazují abstraktním neuronům lidsky pochopitelný význam.

Na tuto transformaci latentního prostoru plynule navazuje mechanismus vysvětlitelnosti, jenž s využitím myšlenek z modelů jako \textit{ReEL (Review-Aware Explanation of Location Recommendation)} \cite{Baral2018} propojuje extrahované faktory s konkrétními tagy a aspekty míst (např. atmosféra, vhodnost pro děti, kvalita služeb). Tím vzniká exaktní základ pro generování personalizovaných vysvětlení. Systém však nezůstává pouze u pasivního vysvětlování. Inspirací z interaktivních nástrojů typu \textit{TasteWeights} \cite{Bostandjiev2012} práce navrhuje mechanismus přímého řízení (\textit{steering}), jímž uživatelé získávají možnost do svých vygenerovaných profilů aktivně zasáhnout. Zviditelnění těchto ovládacích prvků vrací uživatelům nad systémem kontrolu (tzv. \textit{agency}), což dle stávajících výzkumů pomáhá vyvažovat obavy ze ztráty soukromí spojené s profilováním \cite{Knijnenburg2012}. Jak navíc dokazují nedávné studie \cite{Ooge2023}, právě tento přechod od \textit{black-box} doporučení k ovladatelné interakci – za předpokladu, že uživatel přímo vidí okamžitý dopad svých úprav – má prokazatelně signifikantní vliv na budování uživatelské důvěry (trust) a vnímané transparentnosti celého doporučovacího procesu.

\section*{Struktura práce}
Text této diplomové práce je logicky členěn do pěti hlavních kapitol. První kapitola definuje formální a matematický rámec doporučovacích systémů. Představuje architekturu autoenkodérů, principy vynucování řídkosti u skrytých reprezentací a teoreticky ukotvuje pojmy vysvětlitelnosti a interaktivního řízení preferencí (steering). Na tyto základy navazuje druhá kapitola, jež podává ucelenou rešerši současného stavu poznání. Mapuje existující přístupy k vysvětlitelnosti, metody pro rozpletení latentních prostorů (disentanglement) a aktuální trendy ve využívání velkých jazykových modelů v doporučování. 

Třetí kapitola následně představuje samotný architektonický návrh systému. Jsou zde definovány výzkumné otázky a je detailně popsána celá navrhovaná pipeline – od kolaborativního filtrování přes automatické sémantické pojmenovávání extrahovaných neuronů pomocí LLM až po logiku uživatelského rozhraní. Praktické realizaci tohoto návrhu je věnována čtvrtá kapitola. Ta popisuje využité datové sady Yelp, kroky předzpracování, implementačními detaily doporučovacího modelu a fungováním webové aplikace vytvořené pro uživatelské testování. V páté kapitole jsou poté shrnuty a analyzovány dosažené výsledky. Prezentovány jsou nejen objektivní metriky přesnosti na historických datech, ale především detailní vyhodnocení uživatelské studie zaměřené na to, jak koncoví uživatelé vnímají transparentnost a ovladatelnost systému. Závěr práce rekapituluje dosažené výsledky vzhledem k vytyčeným cílům, objektivně diskutuje limity navrženého řešení a nastiňuje možné směry pro budoucí výzkum.
